% Eine Quelle, zwei PDFs. scripts/build.sh setzt \explainmode:
%   0 → loesung_blattNN.pdf     (Abgabe: Angabe + Schritte, schlicht)
%   1 → erklaerung_blattNN.pdf  (identisch + \explain{} / \explainfig{})
\documentclass[a4paper,11pt]{article}
\usepackage[T1]{fontenc}
\usepackage[utf8]{inputenc}
\usepackage[ngerman]{babel}          % englisches Blatt: [english]
\usepackage{amsmath,amssymb,mathtools}
\usepackage[margin=2.2cm]{geometry}
\usepackage{enumitem}
\usepackage{xcolor}
\usepackage{booktabs}
\usepackage{tikz}
\usepackage{pgfplots}\pgfplotsset{compat=1.18}
\usepackage[most]{tcolorbox}
\setlength{\parindent}{0pt}
\setlength{\parskip}{4pt}

% ---- Modus ----
\providecommand{\explainmode}{0}
\newif\ifexplain
\ifnum\explainmode=1 \explaintrue\fi

% ---- Gerüst (in beiden PDFs) ----
\newcommand{\blatttitel}[2]{\begin{center}{\Large\bfseries #1}\\[2pt]#2\end{center}}
\newcommand{\aufgabe}[1]{\section*{Aufgabe #1}}          % englisch: \renewcommand → Exercise
\newcommand{\angabe}[1]{\textit{#1}\par\smallskip}
\newcommand{\teil}[1]{\subsection*{(#1)}}
\newcommand{\why}[1]{\text{\small\color{black!65}#1}}   % Begründung rechts: && \why{...}

% ---- Nur im Erklär-PDF ----
\newtcolorbox{erklbox}{colback=blue!4,colframe=blue!40,breakable,boxrule=0.5pt,
  fonttitle=\bfseries\small,title={Erklärung}}
\newcommand{\explain}[1]{\ifexplain\begin{erklbox}#1\end{erklbox}\fi}
\newcommand{\explainfig}[1]{\ifexplain\begin{center}#1\end{center}\fi}

\begin{document}
\blatttitel{Übungsblatt NN -- \ifexplain Erklärung\else Lösung\fi}{Kursname}

\explain{Kurzer Überblick: Welche Themen kommen vor, welche Begriffe muss man kennen?}

\aufgabe{1}
\angabe{Lösen Sie das Gleichungssystem $\beta = -2$, $0 = \alpha + 2\beta$.}
\explain{Wir setzen den bekannten Wert für $\beta$ in die zweite Gleichung ein. Dann bleibt nur noch eine Unbekannte übrig.}
\begin{align*}
0 &= \alpha + 2\beta    && \why{zweite Gleichung}\\
0 &= \alpha + 2\cdot(-2) && \why{einsetzen: $\beta = -2$}\\
0 &= \alpha - 4          && \why{$2\cdot(-2) = -4$}\\
4 &= \alpha              && \why{$|\,+4$}
\end{align*}
\explain{Das „$|\,+4$“ heißt: Auf beiden Seiten wird $4$ addiert. Links steht dann $4$, rechts bleibt $\alpha$ allein.}
Probe:
\begin{align*}
\alpha + 2\beta &= 4 + 2\cdot(-2) && \why{einsetzen}\\
&= 4 - 4 && \why{$2\cdot(-2)=-4$}\\
&= 0 \;\checkmark && \why{$4 - 4 = 0$}
\end{align*}
\[\boxed{\alpha = 4,\ \beta = -2}\]

\end{document}
